\begin{table}[!tbp]
\centering
\scriptsize
\setlength{\tabcolsep}{1.8pt}
\renewcommand{\arraystretch}{1.12}
\caption{Composition-adjusted comparison of fully recorded (M\_i = 0) vs. known-status cases in partially recorded outbreaks (M\_i $>$ 0, K\_rec $>$ 0), and benchmark rates with 95\% Wilson CIs in 11 low-missingness states (<5\% unrecorded and $\geq$200 common-support cases).}
\label{tab:delta_cal}
\begin{tabular}{>{\raggedright\arraybackslash}p{1.75cm}>{\raggedleft\arraybackslash}p{2.0cm}>{\raggedleft\arraybackslash}p{1.35cm}>{\raggedleft\arraybackslash}p{1.45cm}>{\raggedleft\arraybackslash}p{1.45cm}>{\raggedleft\arraybackslash}p{1.35cm}>{\raggedleft\arraybackslash}p{1.45cm}>{\raggedleft\arraybackslash}p{2.9cm}}
\toprule
Group & Full vs Part (\%) & Within-cell Δ (pp) & Std obs $-$ Full (pp) & Single-state Δ (pp) & Multistate Δ (pp) & 11 states n / C & 11 states rate [95\% CI] \\
\midrule
STEC & 25.3 vs 36.2 & +0.19 & -2.53 & +8.37 & -1.28 & 20 / 264 & 26.2 [21.3,31.9] \\
Salmonella & 14.1 vs 25.9 & +9.02 & +3.21 & +4.34 & +11.98 & 110 / 2,378 & 13.0 [11.7,14.5] \\
Campylobacter & 7.9 vs 10.8 & +6.59 & -0.27 & +2.40 & +34.62 & 28 / 207 & 8.8 [5.5,13.6] \\
Norovirus GII & 1.4 vs 2.1 & +1.79 & +0.13 & +0.72 & -- & 115 / 3,316 & 1.3 [0.9,1.7] \\
C. perfringens & 0.5 vs 0.7 & +1.38 & +0.13 & +0.28 & -- & 18 / 895 & 0.0 [0.0,0.4] \\
Norovirus GI & 0.3 vs 0.8 & +0.05 & +0.19 & +0.44 & -- & 26 / 712 & 0.3 [0.1,1.0] \\
\bottomrule
\end{tabular}
\par\smallskip\parbox{0.97\linewidth}{\footnotesize Within-cell Δ (+0.05 to +9.02 pp) and single-state Δ are positive across all six groups, whereas Std obs $-$ Full is negative for STEC ($-$2.53 pp) and Campylobacter ($-$0.27 pp) and Multistate Δ is negative for STEC ($-$1.28 pp). Common-mechanism scaling (q\_js = $\alpha$ p\_js) preserves 14/15 directions at $\alpha$ = 0 and 15/15 for $\alpha$ $\geq$ 0.50 (Supp. Fig. S4). With a 0.1\% zero-cell floor under OR scaling, 14/15 pairs hold at $\Gamma$ = 3.0 and 13/15 at $\Gamma$ = 5.0.}
\end{table}
